\documentclass[../../../../main.tex]{subfiles}
\begin{document}

\begin{flushright}
\footnotesize\textit{【自動文字起こし・要確認】}
\end{flushright}

[解] $m \ge 2 \cdots (1)$

$d_m$ は $mC_1$ を割り切ることが必要で、すなわち $d_m$ は $m$ の約数である $\cdots (2)$

(1) $m \in \text{prime}$ の時、(2) から $d_m = 1 \text{ or } m$ である。以下、$k=2,\dots,m-1$ に対し、
${}_{m}C_{k}$ が $m$ で割り切れることを示す。
\[
k \cdot {}_{m}C_{k} = m \cdot {}_{m-1}C_{k-1} \quad \therefore \quad {}_{m}C_{k} = \frac{{}_{m-1}C_{k-1}}{k} \cdot m
\]
ここで、$m$ と $k$ は互いに素で、かつ ${}_{m}C_{k} \in \mathbb{N}$ だから、$\frac{{}_{m-1}C_{k-1}}{k} \in \mathbb{N}$ となり、
たしかに ${}_{m}C_{k}$ は $m$ の倍数\fbox{終}

以上から、$d_m = m$ が従う\fbox{終}

(2) $k=1$ の時は成立するから、以下 $k=i \in \mathbb{N}$ での成立を仮定し、$k=i+1$ での成立を示す。$a_m(k) = k^m - k$ とおく。
\begin{align*}
a_m(i+1) &= (i+1)^m - (i+1) = \sum_{t=0}^{m} {}_{m}C_{t} \cdot i^t - i - 1 \\
&= \sum_{t=1}^{m-1} {}_{m}C_{t} \cdot i^t + (i^m - i)
\end{align*}
ここで、${}_{m}C_{t}$、$i^m - i$ は共に $d_m$ で割り切れるので、$a_m(i+1)$ も $d_m$ で割り切れる。
よって $k=i+1$ でも成立し、以上から題意は示された\fbox{終}

(3) $(1-1)^m = \sum_{k=0}^{m} {}_{m}C_{k} (-1)^k = 2 + \sum_{k=1}^{m-1} {}_{m}C_{k} (-1)^k = 0 \quad (\because m \in \text{even}) \cdots (3)$

${}_{m}C_{k} \ (k=1,2,\dots,m-1)$ は $d_m$ で割り切れるので、$S \in \mathbb{Z}$ として、(3) から
\[
d_m \cdot S = 2
\]
とかける。$d_m \in \mathbb{N}$ から、$d_m = 1 \text{ or } 2$ が従う。\fbox{終}

\end{document}
